\chapter{Compatible families and sparse changes}\label{ch:sparse}
\sourcebox{erdos1060\_sparse\_changes\_proof.tex}{Strong restricted-case results. Compatibility is an essential hypothesis; Chapter~\ref{ch:ratio} gives the exact obstruction to deleting it.}
\section{Definitions and scope}
Write
\[
 h(k)=k\sigma(k),\qquad \alpha_p=v_p(N),\qquad
 f_E(N)=\#\{k:h(k)=N,\ v_p(k)\le E\text{ for every prime }p\}.
\]
All logarithms are natural. The assertions involving $w=\log\log N$ are asymptotic as $N\to\infty$, so all iterated logarithms used below are positive. Set
\[
 \PP(N)=\prod_{p\mid N}\frac{p}{p-1},\qquad
 \omega(N)=\#\{p:p\mid N\},\qquad \Omega(N)=\sum_p\alpha_p.
\]
A nonempty subfamily $\FF\subseteq\{k:h(k)=N\}$ is called \emph{compatible} if, for every $k,m\in\FF$ and every prime $p$, one of
\[
 v_p(k)+1,\qquad v_p(m)+1
\]
divides the other. The direction of divisibility may vary with both the pair and the prime. The compatibility hypothesis is essential and is not asserted for an arbitrary fiber.

The new argument records changes in small-prime divisor-sum valuations relative to one reference preimage. This replaces the need to record every block contributing such a valuation inside a compatible family.


\section{The integer multiplier for a compatible pair}
\begin{lemma}\label{sparse:lem:multiplier}
Suppose $k\ne m$, $h(k)=h(m)=N$, and the exponent indices are comparable at every prime. For each differing prime, write the larger exponent as $a_p$, the smaller as $b_p$, and $d_p=a_p-b_p>0$. Then there is an integer $c>1$, dividing $N$, such that
\[
 1<c^2<\PP(N).
\]
If in addition $v_p(k)=v_p(m)$ for all $p\le Y$, then, for each prime $q\le Y$,
\begin{equation}\label{sparse:eq:l1identity}
 \sum_{p>Y}\left|v_q\bigl(\sigma(p^{v_p(k)})\bigr)
                   -v_q\bigl(\sigma(p^{v_p(m)})\bigr)\right|
 =2v_q(c).
\end{equation}
All sums are finite; absent input primes have exponent zero and divisor sum $1$.
\end{lemma}
\begin{proof}
Comparability gives $b_p+1\mid a_p+1$. Thus
\[
 \frac{h(p^{a_p})}{h(p^{b_p})}=p^{d_p}Q_p,
 \qquad Q_p=\frac{p^{a_p+1}-1}{p^{b_p+1}-1}\in\NN,
\]
and the geometric sum gives
\begin{equation}\label{sparse:eq:local}
 1<\frac{Q_p}{p^{d_p}}<\frac1{1-p^{-(b_p+1)}}\le\frac p{p-1}.
\end{equation}
Let $P_+$ contain the primes with larger exponent in $k$, and $P_-$ those with larger exponent in $m$. Put
\[
 A=\prod_{p\in P_+}p^{d_p},\quad B=\prod_{p\in P_-}p^{d_p},\quad
 U=\prod_{p\in P_+}Q_p,\quad V=\prod_{p\in P_-}Q_p.
\]
Cancel equal local factors in $h(k)=h(m)$. Then $AU=BV$ and $\gcd(A,B)=1$, so
\[
 U=cB,\qquad V=cA
\]
for a positive integer $c$. Since $U\mid\sigma(k)$, we have $c\mid N$. Multiplying \eqref{sparse:eq:local} shows
\[
 1<c^2=\frac{UV}{AB}
 <\prod_{p\in P_+\cup P_-}\frac p{p-1}\le\PP(N).
\]
In particular $c>1$.

If all input exponents at primes at most $Y$ agree, no such prime divides $A$ or $B$. Hence, for $q\le Y$,
\[
 v_q(U)=v_q(V)=v_q(c).
\]
At a differing prime $p$, the absolute difference between the two local $q$-adic divisor-sum valuations is $v_q(Q_p)$. Summing over $P_+$ and $P_-$ proves \eqref{sparse:eq:l1identity}.
\end{proof}


\section{Sparse labels determine a compatible preimage}
For $Y\ge2$ and each input base $p>Y$, define the label
\[
 \lambda_p(k)=\bigl(v_q(\sigma(p^{v_p(k)}))\bigr)_{q\le Y,\ q\text{ prime}}.
\]
This is a finite vector of nonnegative integers.

\begin{theorem}\label{sparse:thm:finite}
Let $N\ge2$, let $\FF$ be a compatible subfamily of its fiber, and take
\[
 Y\ge\max\{2,\sqrt{\PP(N)}\},\qquad
 D=\left\lfloor\frac{\log\PP(N)}{\log2}\right\rfloor.
\]
If every input exponent in $\FF$ is at most $E$, then
\begin{equation}\label{sparse:eq:finiteE}
 \boxed{|\FF|\le (E+1)^{\pi(Y)}(1+E\omega(N))^D.}
\end{equation}
Without an input exponent bound,
\begin{equation}\label{sparse:eq:finiteunbounded}
 \boxed{|\FF|\le
 \prod_{p\mid N,\ p\le Y}(\alpha_p+1)\,
 (1+\Omega(N))^D.}
\end{equation}
\end{theorem}
\begin{proof}
First partition $\FF$ according to all exact input exponents at primes $p\le Y$. The number of parts is at most $(E+1)^{\pi(Y)}$ in the bounded case and the displayed product in the unbounded case.

Fix a nonempty part $\mathcal G$. We show that its full vector of large-base labels $(\lambda_p(k))_{p\mid N,\ p>Y}$ is injective. If two different members had the same vector, each integer local quotient $Q_p$ from Lemma~\ref{sparse:lem:multiplier} would have no prime factor at most $Y$. Its integer multiplier $c$, as a divisor of $U$, would have the same property. But
\[
 1<c<\sqrt{\PP(N)}\le Y,
\]
a contradiction. Thus equal label vectors imply equal preimages.

Choose a reference $k_0\in\mathcal G$. For $k\ne k_0$, Lemma~\ref{sparse:lem:multiplier} gives
\[
 \sum_{p>Y}\sum_{q\le Y}
 \left|v_q(\sigma(p^{v_p(k)}))-v_q(\sigma(p^{v_p(k_0)}))\right|
 =2\sum_{q\le Y}v_q(c)\le2\Omega(c).
\]
Each changed label contributes at least one to the left side. Also
\[
 2\Omega(c)\le\frac{2\log c}{\log2}
 <\frac{\log\PP(N)}{\log2}.
\]
Consequently at most $D$ labels can differ from those of $k_0$. The assertion is immediate for $k=k_0$ as well.

For a fixed base $p$, at most $E$ label values differ from its reference value, since at most $E+1$ exponents are allowed. Encode the changed bases and their replacement labels as a sorted list, padded with empty entries to length $D$. Its alphabet has at most $1+E\omega(N)$ elements. The injectivity already proved therefore gives
\[
 |\mathcal G|\le(1+E\omega(N))^D.
\]
In the unbounded case there are at most $\alpha_p$ replacement labels at $p$, so the alphabet size is at most $1+\sum_p\alpha_p=1+\Omega(N)$. Summing over the parts proves both bounds.
\end{proof}


\section{Uniform asymptotic consequences}
\begin{lemma}\label{sparse:lem:mertens}
Uniformly over integers $N$ tending to infinity,
\[
 \PP(N)\ll\log\log N.
\]
\end{lemma}
\begin{proof}
Put $L=\log N$ and $w=\log L$. The part with primes at most $L$ is at most
\[
 \prod_{p\le L}\frac p{p-1}\ll\log L=w
\]
by Mertens' product theorem \cite{mertens}. The number of prime divisors of $N$ larger than $L$ is at most $L/\log L$. Their logarithmic contribution is at most
\[
 \sum_{p\mid N,\ p>L}\frac1{p-1}=O(1/w).
\]
Multiplying the two estimates proves the lemma.
\end{proof}

\begin{corollary}\label{sparse:cor:compatible}
Let $w=\log\log N$, and let $\FF$ be a compatible subfamily of the fiber of $N$.
For each fixed $E$, if all input exponents are at most $E$, then
\begin{equation}\label{sparse:eq:compatgrowth}
 \log|\FF|=O_E(w\log w).
\end{equation}
Without any input exponent bound,
\[
 \log|\FF|=O(w^{3/2}).
\]
Both are $o(\log N/\log\log N)$.
\end{corollary}
\begin{proof}
Choose $Y=\max\{2,\sqrt{\PP(N)}\}$. Lemma~\ref{sparse:lem:mertens} gives
\[
 Y=O(\sqrt w),\qquad D=O(\log w),\qquad
 \log(1+\Omega(N))=O(w).
\]
Since $\pi(Y)\le Y$ and $\omega(N)\le\Omega(N)\le\log N/\log2$, \eqref{sparse:eq:finiteE} gives \eqref{sparse:eq:compatgrowth}. For \eqref{sparse:eq:finiteunbounded}, the logarithm of its small-prime factor is at most
\[
 \pi(Y)\log(1+\log N/\log2)=O(w^{3/2}).
\]
The remaining term is $O(w\log w)$, which is smaller asymptotically. The final claim follows because $\log N/\log\log N=e^w/w$.
\end{proof}

\begin{corollary}\label{sparse:cor:boundedtarget}
Fix $E,A$. If $\FF$ is compatible, every input exponent is at most $E$, and $\alpha_p\le A$ for every prime $p$, then
\begin{equation}\label{sparse:eq:boundedtarget}
 \log|\FF|=O_{E,A}\!\left(\frac{w\log w}{\log\log w}\right).
\end{equation}
\end{corollary}
\begin{proof}
The multiplier $c$ in Lemma~\ref{sparse:lem:multiplier} divides $N$, and $c<\sqrt{\PP(N)}=O(\sqrt w)$. If $s=\omega(c)$, then
\[
 c\ge\prod_{j=1}^{s}p_j\ge(s+1)!,\qquad \Omega(c)\le As.
\]
The elementary factorial bound gives
\[
 s=O\!\left(\frac{\log w}{\log\log w}\right).
\]
Indeed $\log c=O(\log w)$, while $\log(s!)\ge(s/2)\log(s/2)$ for $s\ge4$; bounded $s$ is harmless. Therefore the number of changed labels in the proof of Theorem~\ref{sparse:thm:finite} is at most
\[
 O_A\!\left(\frac{\log w}{\log\log w}\right)
\]
in place of $D$. Each replacement record costs $O_E(w)$ logarithmically, and the small-prime exponent record costs $O_E(\sqrt w)$. This proves \eqref{sparse:eq:boundedtarget}.
\end{proof}


\section{Reducing bounded-exponent fibers to compatible parts}
\begin{theorem}\label{sparse:thm:cap3}
Uniformly for targets $N\to\infty$, with $w=\log\log N$,
\begin{equation}\label{sparse:eq:cap3}
 \boxed{\log\max(1,f_3(N))
 =O\bigl((v_3(N)+\log w)w\bigr).}
\end{equation}
\end{theorem}
\begin{proof}
Fix the exact exponent at the input prime $3$, with at most four choices. Among all other positive input blocks $p^e$, record those for which $3\mid\sigma(p^e)$. There are at most $v_3(N)$ such blocks, because each contributes at least one to $v_3(\sigma(k))\le v_3(N)$. Their bases divide $N$ and their exponents are at most $3$. Thus the number of records is at most
\[
 4(1+3\omega(N))^{v_3(N)}.
\]
As before, this follows by using sorted lists padded with empty entries.

Fix a record and remove its full input prime-power blocks. Every remaining base $p\ne3$ satisfies $3\nmid\sigma(p^e)$. If $p\equiv1\pmod3$, the possible exponent indices for $0\le e\le3$ are $1,2,4$. If $p\equiv-1\pmod3$, they are $1,3$. In either case they form a divisibility chain. At recorded bases the exponents are fixed. Hence the entire part is compatible.

Applying Corollary~\ref{sparse:cor:compatible} with $E=3$ to each part and counting the records proves \eqref{sparse:eq:cap3}. More explicitly, with $Y,D$ as in Theorem~\ref{sparse:thm:finite},
\[
 f_3(N)\le4^{1+\pi(Y)}(1+3\omega(N))^{v_3(N)+D}.
\]
\end{proof}

\begin{theorem}\label{sparse:thm:cap5}
Uniformly for targets $N\to\infty$,
\begin{equation}\label{sparse:eq:cap5}
 \boxed{\log\max(1,f_5(N))
 =O\bigl((v_2(N)+v_3(N)+v_5(N)+\log w)w\bigr).}
\end{equation}
\end{theorem}
\begin{proof}
Fix the exact input exponents at $2,3,5$, with at most $6^3$ choices. For other positive input blocks, record those whose divisor sum is divisible by $2$, $3$, or $5$. There are at most
\[
 b=v_2(N)+v_3(N)+v_5(N)
\]
such blocks, so there are at most $6^3(1+5\omega(N))^b$ records.

Fix a record, let $d$ be the product of all recorded full prime-power blocks, and write $k=du$ with $\gcd(d,u)=1$. The residuals all satisfy $h(u)=N/h(d)$. Their prime bases exceed $5$ and $\gcd(\sigma(u),30)=1$. Consequently all positive exponents are $2$ or $4$, because an odd input base to an odd exponent has an even divisor sum.

We show that any two residuals in this part have comparable indices at each prime. The only possible incomparable pair is $3,5$, arising from exponents $2,4$ at a base $p>5$. Coprimality of the corresponding divisor sums with $3$ and $5$ forces
\[
 p\not\equiv1\pmod3,\qquad p\not\equiv1\pmod5.
\]
On the other hand, every prime $q>5$ dividing a local divisor sum with index $3$ or $5$ is $1$ modulo $3$ or $5$, respectively. To check this, $q$ divides $x^\ell-1$ with $\ell=3$ or $5$; if $x\equiv1\pmod q$ then $q\mid\ell$, impossible, and otherwise the order of $x$ modulo $q$ is $\ell$. Thus such a $p$ divides neither residual divisor sum. Taking its valuation in the equality of their $h$-values forces equal input exponents, a contradiction.

The residuals, and hence the original part, are compatible. Corollary~\ref{sparse:cor:compatible} proves \eqref{sparse:eq:cap5}. The finite bound is
\[
 f_5(N)\le6^{3+\pi(Y)}(1+5\omega(N))^{b+D}.
\]
\end{proof}


\section{Two full-fiber consequences}
\begin{theorem}\label{sparse:thm:unboundedtwo}
Let $N=2^a m$, where $m$ is odd and fourth-power-free; equivalently, $v_p(m)\le3$ for every prime $p$. There is no restriction on $a$. Uniformly over these targets,
\[
 \boxed{\log\max(1,f(N))=O(w\log w).}
\]
In particular the zero-constant bound of Erd\H{o}s Problem 1060 holds on this class.
\end{theorem}
\begin{proof}
Fix the exponent $e=v_2(k)$ of a preimage, at a cost of at most $a+1$ choices, and write $k=2^e u$ with $u$ odd. The residual target $M=N/h(2^e)$ is integral for an actual preimage, and $h(u)=M$. All input exponents in $u$ are at most $3$, since $u\mid m$, and $v_3(M)\le v_3(N)\le3$.

Theorem~\ref{sparse:thm:cap3}, applied to $M\le N$, bounds each nonempty residual fiber by $\exp(O(w\log w))$. The estimate can be used uniformly up to $N$ by enlarging its absolute constant to include bounded $M$. Finally,
\[
 \log(a+1)\le\log(1+\log N/\log2)=O(w).
\]
Multiplying by the number of choices for $e$ proves the result.
\end{proof}

\begin{theorem}\label{sparse:thm:sixth}
Uniformly over sixth-power-free targets $N$, meaning $v_p(N)\le5$ for all primes $p$,
\[
 \boxed{\log\max(1,f(N))
 =O\!\left(\frac{w\log w}{\log\log w}\right).}
\]
This improves the previous restricted-case bound $O((w/\log w)^{3/2})$ in the earlier working note.
\end{theorem}
\begin{proof}
Every preimage divides $N$, so $f(N)=f_5(N)$. In the partition from Theorem~\ref{sparse:thm:cap5}, we have $b\le15$, and hence only $\exp(O(w))$ records. Each part is compatible. Apply Corollary~\ref{sparse:cor:boundedtarget} with $E=A=5$ to its size. The cost $O(w)$ of the records is smaller than the stated bound.
\end{proof}

For either theorem, its logarithmic upper bound divided by $\log N/\log\log N=e^w/w$ tends to zero. These results concern the full multiplicity on their stated target classes, not merely a bounded-exponent subset of the fiber.

